\makeatletter
% Add the ICLR 2026 template directory to LaTeX's search path.
% This lets \usepackage{iclr2026_conference} work without copying files.
\def\input@path{{../doc/iclr-template/master/iclr2026/}}
\makeatother

\documentclass{article}

% ICLR template (located in ../doc/templates/iclr2026/)
\usepackage{iclr2026_conference}
\usepackage{times}

% Math & symbols
\usepackage{amsmath,amsfonts,amssymb,amsthm}

% Figures & tables
\usepackage{graphicx}
\usepackage{booktabs}
\usepackage{subcaption}

% Formatting
\usepackage{microtype}
\usepackage{hyperref}
\usepackage{xcolor}

% Reduce underfull vbox warnings by allowing variable page height.
\raggedbottom

% Theorem-like environments (used in \input files).
\newtheorem{theorem}{Theorem}
\newtheorem{lemma}{Lemma}
\newtheorem{proposition}{Proposition}
\newtheorem{corollary}{Corollary}
\newtheorem{assumption}{Assumption}

\title{Beyond the Hessian: Jacobian Anisotropy and Learning Dynamics in Large Language Models}

\author{
  Anonymous Authors\\
  Anonymous Institutions\\
  \texttt{anonymous@iclr.cc}
}

% Uncomment for camera-ready
% \iclrfinalcopy

\begin{document}

\maketitle

\begin{abstract}

Classical learning theory attributes the stability and convergence of
gradient descent to curvature, with learning rate constraints determined
by the largest eigenvalue of the Hessian.
Recent analyses of deep networks show that training often operates near
the edge of stability, yet these works treat the underlying curvature
structure as given.
In large language models (LLMs), the autoregressive next-token objective,
cross-entropy loss, and self-attention architecture induce persistent
statistical anisotropy in learned representations.
From first principles, language modeling produces nonzero mean components
and low-rank covariance spikes, yielding structured spectral concentration
in the Jacobian.

In this work, we develop a learning-theoretic framework for gradient descent
under Jacobian anisotropy motivated by large language models.
We assume a spike--tail decomposition of the Jacobian and derive
stability conditions, maximal admissible learning rates,
and convergence rates that explicitly depend on spike magnitude.
Our analysis shows that learning dynamics are governed by low-rank
spectral structure rather than by abstract smoothness constants alone.
Furthermore, we formalize how normalization, mean removal,
and spectral whitening reduce dominant singular components,
thereby relaxing learning-rate constraints and accelerating
tail convergence.

Together, our results provide a Jacobian-centric theory of learning
dynamics that extends classical Hessian-based analysis to the
domain-specific geometry of large language model training.

\end{abstract}

\section{Introduction}

Large language models are trained on massive human language corpora,
which induce a characteristic geometry in learned representations.
Across layers and training stages, activations, gradients, and model parameters
exhibit pronounced anisotropy:
a small number of directions concentrate disproportionate energy
(strong spectral spikes),
while the remaining directions form a broad long tail
\cite{ethayarajh2019contextual,mu2018all,timkey2021rogue,kovaleva2021bertbusters,machina2024anisotropy}.

This anisotropy is not accidental.
It reflects the hierarchical structure of human language itself.
At every level of representation—tokens, phrases, clauses, and full sentences—
there exist shared semantic components that are common across tokens within a context.
When contextual aggregation mechanisms such as attention and nonlinear transformations
extract these shared semantics, they naturally induce coherent mean directions
in the representation space.
Consequently, mean bias is not confined to the embedding layer;
it re-emerges across layers as models progressively capture higher-level
semantic regularities.

While such structure is linguistically natural,
it becomes a computational liability during large-scale training.
Anisotropy inflates elementwise dynamic range,
distorts adaptive second-moment estimates
\cite{kingma2014adam,loshchilov2017adamw,spectra2026},
and compresses long-tail variation in low-bit regimes
\cite{xiao2023smoothquant,frantar2022gptq,lin2024awq,li2024svdquant,ashkboos2024quarot}.
In high-dimensional spaces, even moderate coherent bias produces
a consistent directional drift across thousands of coordinates,
amplifying outliers and destabilizing numerical precision.
Thus anisotropy constitutes a structural \emph{curse}
for efficient and stable large-scale training.

In this work, we perform a systematic layer-wise investigation
of anisotropy across activations, gradient outputs, and model parameters.
We uncover a structural asymmetry:
\emph{mean bias is a dominant component of activation anisotropy across layers,
whereas gradient outputs exhibit no comparable coherent mean}.
Moreover, although model parameters display spectral spikes,
they do not accumulate persistent mean offsets.
This establishes activation mean bias as the principal forward-path driver
of anisotropic geometry.

Importantly, we show that anisotropy can be decomposed into two components:
(i) a coherent mean bias, and
(ii) a residual low-rank spike.
The mean bias is the primary source of extreme elementwise outliers
and adaptive second-moment inflation,
while the residual spike plays a secondary role.
Thus the dominant instability mechanism is not generic high-dimensional variance,
but a structured rank-one component that propagates through residual connections
across depth.

Tracing this structure backward,
we identify its statistical origin in the skewed token distribution of natural language.
Embedding weights develop frequency-aligned structure
\cite{mikolov2013efficient,mikolov2013distributed,arora2017simple},
and subsequent contextual aggregation amplifies shared semantics at progressively
higher abstraction levels.
Mean bias therefore reflects the hierarchical organization of world knowledge
encoded in language rather than an architectural accident.

Crucially, this structure is not merely harmful—it is also exploitable.
Because the dominant instability arises from a rank-one component,
it can be isolated and removed with a single mean-subtraction operation.
The remaining low-rank spike can, if necessary,
be whitened once at the embedding activation level.
Both operations are hardware-friendly and substantially cheaper
than repeated layer-wise SVD corrections.
In this sense, the same structure that creates instability
also enables efficient mitigation—the \emph{blessing} of bias.

We validate this bias-centric interpretation in two downstream settings.
In low-bit quantization, mean removal significantly reduces dynamic-range inflation
\cite{xiao2023smoothquant,wei2023outlier,ashkboos2024quarot}.
In adaptive optimization, it contracts extreme second-moment tails,
improving effective learning-rate allocation
\cite{kingma2014adam,spectra2026}.
These results suggest that addressing bias at its source
can replace or complement spectral interventions
with minimal computational overhead.

\paragraph{Contributions.}
\begin{enumerate}[leftmargin=*,itemsep=2pt]
    \item We show that activation mean bias emerges across layers as a structural consequence of hierarchical semantic aggregation.
    \item We demonstrate that mean-dominated components account for most extreme activation magnitudes and second-moment inflation.
    \item We establish a forward-path asymmetry: gradient outputs remain mean-neutral while activations accumulate coherent bias.
    \item We propose source-level mean removal and optional spike whitening as efficient spectral conditioning mechanisms.
    \item We demonstrate improvements in low-bit quantization and adaptive optimization.
\end{enumerate}
\section{Related Work}

\paragraph{Learning-rate stability and the Hessian paradigm.}
Classical convergence and stability analyses of gradient descent are typically
stated in terms of smoothness constants or Hessian curvature, leading to
step-size constraints proportional to $1/\lambda_{\max}(H)$.
This Hessian-centric viewpoint underlies much of modern optimization theory
for machine learning and motivates the common practice of diagnosing training
instabilities through curvature (sharpness) measurements.

\paragraph{Edge of Stability (EoS) and oscillatory optimization dynamics.}
Cohen et al.~\cite{cohen2021eos} empirically identified that full-batch gradient descent
in neural network training often operates in an \emph{Edge-of-Stability} regime,
where the top Hessian eigenvalue fluctuates near $2/\eta$ and the loss exhibits
short-term oscillations while decreasing over long horizons.
Arora et al.~\cite{arora2022eos} provided a theoretical analysis of mechanisms by which
gradient-based methods can enter and make progress in such EoS regimes,
thereby initiating a line of work on understanding learning dynamics beyond
the classical monotone-descent picture.
More recently, Cohen et al.~\cite{cohen2024centralflows} proposed \emph{central flows},
a differential-equation formalism aimed at modeling the time-averaged trajectory
of oscillatory optimizers in EoS-like regimes, providing a complementary tool
for reasoning about long-term dynamics.

\paragraph{Beyond GD: EoS-style analyses for alternative update rules.}
Extensions of the EoS perspective have also appeared for algorithms that modify
the objective or update rule.
For example, Long and Bartlett~\cite{long2024sameos} derive an ``edge of stability''
condition for Sharpness-Aware Minimization (SAM) and empirically validate that
SAM operates near this boundary; see also Foret et al.~\cite{foret2020sam} for
the original SAM formulation.

\paragraph{Matrix-aware preconditioning and matrix-whitening optimizers.}
A separate but related thread studies \emph{structure-aware} optimizers that exploit
the matrix structure of gradients rather than treating parameters as vectors.
PolarGrad~\cite{lau2025polargrad} provides a unifying perspective for ``matrix-aware''
preconditioning and frames Muon as a special instance (via polar decomposition / related
matrix transforms), with both theoretical and empirical evidence of acceleration.
In parallel, practical scaling work has shown that Muon can be made viable for large-scale
language-model training through techniques such as weight decay and per-parameter scaling~\cite{liu2025muonscalable}.
At the same time, the rapidly growing space of matrix-whitening methods has motivated
careful ablations; Frans et al.~\cite{frans2025matrixwhitening} argue that performance
differences among whitening-style optimizers are not explained solely by the accuracy
of spectral normalization, highlighting the role of variance adaptation.

\paragraph{Spike--tail spectral geometry and spike-aware spectral shaping.}
Motivated by spectral anisotropy in large language model training, Huang et al.~\cite{huang2026spectra}
propose Spectra, a spike-aware optimizer that tracks and suppresses a dominant low-rank spike
subspace without amplifying the spectral tail.
Our work is complementary in spirit to these optimizer developments but is primarily
\emph{learning-theoretic}: we focus on deriving stability and convergence properties of
gradient dynamics under a domain-motivated spike--tail Jacobian structure (including mean bias),
and then interpret normalization, mean-removal, and whitening-style interventions
as geometric operations that reduce dominant singular components and relax step-size constraints.
\section{Mean Removal and What the Objective Actually Learns}
\label{sec:mean_removal_from_loss}

This section answers a foundational question: if the token-mean of activations is a random variable that depends on context, why is it nevertheless plausible that removing (or separating) the mean does not destroy the learning signal?
We proceed from the training objective of LLMs, derive the exact gradient expressions, and establish which parts of the signal correspond to the mean (first moment) and which correspond to token-wise variability (second-order structure).

\subsection{Training Objective and Notation}

Large language models are trained by minimizing next-token cross-entropy.
Let $\mathcal{D}$ denote the data distribution over sequences, and let the model define conditional probabilities
\begin{equation}
p_\theta(w_{t+1}\mid w_{\le t}) = \mathrm{softmax}(z_t)_ {w_{t+1}},
\end{equation}
where $z_t\in\mathbb{R}^V$ are logits over vocabulary size $V$.
The population loss is
\begin{equation}
\label{eq:ce_loss}
\mathcal{L}(\theta)=\mathbb{E}_{(w_1,\dots,w_T)\sim\mathcal{D}}\Bigg[-\sum_{t=1}^{T-1}\log \mathrm{softmax}(z_t)_{w_{t+1}}\Bigg].
\end{equation}

Fix a layer $\ell$ and focus on a single linear projection within that layer (the argument applies to attention and MLP projections alike).
Let $x_t\in\mathbb{R}^H$ be the input activation at token $t$ to this projection, and define the projection
\begin{equation}
\label{eq:linear_proj}
y_t = W x_t + b,\qquad W\in\mathbb{R}^{d\times H},\; b\in\mathbb{R}^d.
\end{equation}
Downstream nonlinearities and further layers map $\{y_t\}$ to logits $\{z_t\}$.
Let $\delta_t\in\mathbb{R}^d$ denote the backpropagated gradient at this projection:
\begin{equation}
\label{eq:delta_def}
\delta_t \;\triangleq\; \frac{\partial \mathcal{L}}{\partial y_t}.
\end{equation}
Then the per-token gradient of $\mathcal{L}$ with respect to $W$ is exactly
\begin{equation}
\label{eq:grad_outer}
\nabla_W \mathcal{L} \;=\; \sum_{t} \delta_t x_t^\top,
\qquad
\nabla_b \mathcal{L} \;=\; \sum_t \delta_t.
\end{equation}
This is the precise ``outer-product'' structure used throughout optimization.

\subsection{Mean Decomposition and Exact Gradient Split}

Let $N$ denote the number of tokens contributing to the gradient in a batch (or in an expectation over data).
Define the \emph{token-mean activation} and \emph{centered activation}:
\begin{equation}
\label{eq:mean_center_def}
\mu_x \;\triangleq\; \frac{1}{N}\sum_{t=1}^N x_t,
\qquad
\tilde{x}_t \;\triangleq\; x_t - \mu_x,
\qquad
\sum_{t=1}^N \tilde{x}_t = 0.
\end{equation}
Substitute $x_t=\mu_x+\tilde{x}_t$ into \eqref{eq:grad_outer}. This yields an \emph{exact identity}:
\begin{equation}
\label{eq:grad_split}
\nabla_W \mathcal{L}
=
\sum_t \delta_t(\mu_x+\tilde{x}_t)^\top
=
\Big(\sum_t \delta_t\Big)\mu_x^\top
+
\sum_t \delta_t \tilde{x}_t^\top.
\end{equation}
Define the mean of backprop signals
\begin{equation}
\label{eq:mu_delta_def}
\mu_\delta \;\triangleq\; \frac{1}{N}\sum_{t=1}^N \delta_t.
\end{equation}
Then \eqref{eq:grad_split} can be written as
\begin{equation}
\label{eq:grad_split_compact}
\frac{1}{N}\nabla_W \mathcal{L}
=
\underbrace{\mu_\delta \mu_x^\top}_{\text{rank-1 mean--mean term}}
+
\underbrace{\frac{1}{N}\sum_t \delta_t\tilde{x}_t^\top}_{\text{cross-covariance term}}.
\end{equation}

\paragraph{Interpretation.}
Equation~\eqref{eq:grad_split_compact} is the end-to-end decomposition that matters for learning.
The first term is rank-one and depends only on the \emph{means} $(\mu_\delta,\mu_x)$.
The second term depends on token-wise \emph{variations} $\tilde{x}_t$ and is a cross-covariance between
backprop signals and token-specific deviations.
If the optimization objective requires learning fine-grained distinctions between tokens and contexts, that information must appear in the cross-covariance term.

\subsection{A Fully Rigorous Invariance Statement: Centering with Bias Compensation}

Mean removal in the forward pass changes the function \emph{unless} we compensate for the removed mean through an equivalent parameter transformation.
We now state and prove the exact equivalence.

\begin{theorem}[Exact functional equivalence under mean-centering with bias update]
\label{thm:centering_equiv}
Fix a batch (or context) producing activations $\{x_t\}_{t=1}^N$ and mean $\mu_x$.
Define centered activations $\tilde{x}_t = x_t - \mu_x$.
Consider the linear map \eqref{eq:linear_proj}.
Define an adjusted bias
\begin{equation}
\label{eq:bprime}
b' \;\triangleq\; b + W\mu_x.
\end{equation}
Then for every token $t$,
\begin{equation}
\label{eq:exact_equiv}
W x_t + b \;=\; W\tilde{x}_t + b'.
\end{equation}
Consequently, replacing $x_t$ by $\tilde{x}_t$ \emph{and} replacing $b$ by $b'$ leaves all downstream activations and logits unchanged, hence leaves the cross-entropy loss \eqref{eq:ce_loss} unchanged.
\end{theorem}

\begin{proof}
By definition, $x_t=\tilde{x}_t+\mu_x$. Substitute into the left-hand side:
\[
W x_t + b = W(\tilde{x}_t+\mu_x)+b = W\tilde{x}_t + (W\mu_x+b).
\]
With $b'$ defined by \eqref{eq:bprime}, we obtain \eqref{eq:exact_equiv}.
Since $y_t$ is unchanged for all $t$, every downstream deterministic computation produces identical logits $z_t$, and therefore the softmax probabilities and the loss \eqref{eq:ce_loss} are identical.
\end{proof}

\paragraph{Consequence.}
Theorem~\ref{thm:centering_equiv} is the precise sense in which ``mean removal need not destroy predictive information'':
it is a \emph{reparameterization} of the same function, provided the mean is forwarded (absorbed) as an equivalent bias contribution at the next linear boundary.
This theorem is exact and does not assume anything about the distribution of $\mu_x$.

\subsection{What If We Center Without Compensation?}

If we replace $x_t$ by $\tilde{x}_t$ but keep the original bias $b$, then the linear output becomes
\[
y_t' = W\tilde{x}_t + b = y_t - W\mu_x,
\]
which is a constant shift in $y_t$ across tokens for the given batch.
This shift generally changes logits and therefore changes the loss.
Thus ``centering is free'' is not true without qualification.
The correct statement is: centering is information-preserving conditional on forwarding/absorbing the mean into a separate channel (e.g., adjusted bias, residual carrier, or explicit structure stream).

\subsection{Why ``Useful Information Lives in Variability''}

Equation~\eqref{eq:grad_split_compact} now makes the key point precise.
The only part of the gradient that can encode token- and context-specific distinctions is the cross-covariance term:
\begin{equation}
\label{eq:cross_cov}
C_{\delta x}\;\triangleq\;\frac{1}{N}\sum_{t=1}^N \delta_t\tilde{x}_t^\top.
\end{equation}
The rank-one mean--mean term $\mu_\delta\mu_x^\top$ is shared across tokens and cannot express within-batch variation.

In particular, if $\mu_\delta$ is small (noise-level)---as observed empirically for gradient outputs in our measurements---then the mean--mean term is negligible and learning is driven almost entirely by $C_{\delta x}$.
Even when $\mu_\delta$ is not identically zero, the rank-one term represents a global shared update direction, whereas $C_{\delta x}$ captures structured correlations between errors and token-level deviations, which is the mechanism required to learn conditional dependencies beyond base rates.

\subsection{Summary}

Starting from the cross-entropy loss \eqref{eq:ce_loss}, we derived an exact gradient decomposition \eqref{eq:grad_split_compact} showing that updates split into a rank-one mean--mean component and a cross-covariance component.
We then proved an exact invariance theorem (Theorem~\ref{thm:centering_equiv}) showing that mean-centering can be made function-preserving when the removed mean is forwarded as an equivalent bias at the next linear boundary.
Therefore, mean separation is not inherently information-destroying; information is lost only if the mean channel is discarded rather than separated.



\section{Overview of Key Theoretical Themes}

This paper develops a learning-theoretic framework for gradient descent
under domain-induced Jacobian anisotropy motivated by large language
model (LLM) training. The key themes of the paper are summarized below.

\subsection*{(1) Domain-Specific Structural Prior}

Large language models are trained under an autoregressive next-token
prediction objective with cross-entropy loss. The statistical structure
of human language, together with self-attention architectures, induces
persistent anisotropy in learned representations. From first principles,
we model this structure through a spike--tail decomposition of the
Jacobian:
\begin{equation}
J = J_{\text{mean}} + J_{\text{spike}} + J_{\text{tail}},
\end{equation}
where $J_{\text{mean}}$ captures mean-induced bias, $J_{\text{spike}}$
captures low-rank covariance concentration, and $J_{\text{tail}}$
represents long-tail semantic variation.

\subsection*{(2) Jacobian-Centric Learning Rate Theory}

Under the Gauss--Newton approximation,
\begin{equation}
H \approx J^\top \Lambda J,
\end{equation}
the dominant curvature structure is inherited from the singular spectrum
of $J$. We show that when $\|J_{\text{mean}} + J_{\text{spike}}\|_2
\gg \|J_{\text{tail}}\|_2$, the maximal admissible learning rate is governed
by the dominant singular component:
\begin{equation}
\eta_{\max} \sim \frac{2}{\sigma_{\max}^2(J_{\text{mean}} + J_{\text{spike}})}.
\end{equation}
Thus, learning-rate stability is controlled by low-rank Jacobian geometry
rather than by abstract smoothness constants alone.

\subsection*{(3) Hierarchical Convergence Dynamics}

Spike--tail spectral concentration produces asymmetric learning dynamics.
While the spike direction determines stability,
convergence in tail subspaces is constrained by the same learning-rate bound.
We characterize the convergence rates in spike and tail directions and
quantify the suppression of long-tail semantic components under spike-dominated
regimes.

\subsection*{(4) Edge-of-Stability as a Structural Consequence}

Operating near the boundary
\begin{equation}
\eta \sigma_{\max}^2(J) \approx 2
\end{equation}
naturally produces oscillatory dynamics in the dominant subspace while
maintaining monotone decay in tail directions.
Within our framework, edge-of-stability behavior emerges as a structural
consequence of domain-induced anisotropy.

\subsection*{(5) Spectral Interventions and Geometric Regularization}

Normalization, mean removal, and spectral whitening can be interpreted
as geometric operations that attenuate dominant singular components.
By reducing $\sigma_{\max}(J)$, these interventions enlarge the admissible
learning-rate region and accelerate tail convergence.
We formalize this effect and derive corresponding improvements in
stability and convergence bounds.

\subsection*{(6) A Jacobian-Centric Extension of Classical Theory}

Taken together, our results extend the classical Hessian-based learning
rate paradigm by incorporating domain-specific anisotropy into the
analysis of gradient descent.
The central theoretical claim of this work is that learning dynamics
in large language models are governed by structured Jacobian geometry,
with mean bias and low-rank spikes determining stability, convergence,
and effective learning rates.

\newpage

\bibliographystyle{../doc/iclr-template/master/iclr2026/iclr2026_conference}
\bibliography{references}


\newpage

\appendix


%\newpage
%\input{./sections/learning}


%\input{./draft}

\end{document}